% Reviewed source SHA256 (UTF-8/LF): c648f71278eb5b60b59a9e5fa837ecd623e25d918a05ff3ad54ad245ae4f0eee
\documentclass[11pt]{article}
\usepackage[T1]{fontenc}
\usepackage{lmodern}
\usepackage[a4paper,margin=27mm]{geometry}
\usepackage{amsmath,amssymb,amsthm,mathtools}
\usepackage{microtype,booktabs,array}
\usepackage[hidelinks]{hyperref}
\usepackage{enumitem}
\setlist{itemsep=3pt,topsep=5pt}
\newtheorem{theorem}{Theorem}[section]
\newtheorem{lemma}[theorem]{Lemma}
\newtheorem{corollary}[theorem]{Corollary}
\newtheorem{proposition}[theorem]{Proposition}
\theoremstyle{definition}
\newtheorem{definition}[theorem]{Definition}
\newtheorem{remark}[theorem]{Remark}
\newcommand{\R}{\mathbb R}
\newcommand{\E}{\mathbb E}
\newcommand{\Prob}{\mathbb P}
\newcommand{\tr}{\operatorname{tr}}
\newcommand{\rank}{\operatorname{rank}}
\newcommand{\diag}{\operatorname{diag}}
\newcommand{\range}{\operatorname{range}}
\newcommand{\OPT}{\operatorname{OPT}}
\newcommand{\norm}[1]{\left\lVert#1\right\rVert}
\newcommand{\ip}[2]{\left\langle#1,#2\right\rangle}
\newcommand{\psd}{\succeq}
\newcommand{\pd}{\succ}
\newcommand{\eps}{\varepsilon}
\DeclareMathOperator*{\argmin}{arg\,min}
\allowdisplaybreaks[1]
\hypersetup{pdfcreator={LaTeX},pdfauthor={Matthew J. Colbrook}}

\title{A lower-tail proof for pure relative structured matrix learning}
\author{Matthew J. Colbrook\\\small Department of Applied Mathematics and Theoretical Physics\\\small University of Cambridge, Cambridge, United Kingdom\\\small\href{mailto:m.colbrook@damtp.cam.ac.uk}{m.colbrook@damtp.cam.ac.uk}}
\date{11 September 2026}
\usepackage{xurl}
\setlength{\emergencystretch}{3em}
\begin{document}
\maketitle

\noindent\textbf{Verified affirmative resolution of RE-05.} Theorem 2.1 and Proposition 5.1 give the required query guarantee and confidence. The auxiliary convex extension includes the independently checked failure guard.
Independent agent verification is not external human peer review or formal certification. Original draft remarks about review or source access reflect the input date; the dated \href{https://github.com/MColbrook/OpenProblemsInNLA/blob/codex/colbrook-transfer-resolutions/references/colbrook-transfer-2026-09-11/verification/reviews/RE-05-review.md}{independent review} records the subsequent checks.
\par\medskip
% BEGIN REVIEWED BODY
\begin{abstract}
We give a self-contained proof that a known $q$-dimensional linear matrix family can be fitted to an unknown matrix with pure relative Frobenius error using nonadaptive, two-sided matrix--vector queries. The algorithm measures a high-leverage left subspace exactly and uses a Gaussian right sketch on its orthogonal complement. Only a lower bound on the sketched Gram matrix is required. A second-moment argument and an elementary matrix lower-tail estimate give a constant-success query bound of order $\sqrt{q(\log q+1/\eps)}+\log q$, with no additive error and no dependence on the ratio of the target norm to the optimal residual. A median-radius selector amplifies confidence for affine families. The constant-success argument also extends to closed convex subsets of a known affine family, subject to solving the sketched convex optimization problem.
\end{abstract}

\section{Prior work and precise status}
Amsel et al.\ study structured approximation from two-sided matrix--vector queries and leave a pure relative-error question for linear families \cite{Amsel2026}. A public manuscript repository under the account \texttt{andotheror} subsequently announced the same spectral-splitting algorithm and the pure relative-error conclusion \cite{Prior2026}. Its displayed first-publication filenames are dated July~30, 2026. The proof below was derived independently, but \emph{no priority is claimed for the algorithm or for resolving the pure relative-error question}. The displayed prior query profile contains a squared logarithm. The argument here requires only a lower-tail estimate and gives a single logarithm in that term; comparison with the complete prior manuscript is still necessary before claiming a distinct improvement. The full prior manuscript and its deposit metadata were not available for that comparison.

\section{Model, algorithm, and guarantee}
Let $\mathcal L\subseteq\R^{m\times n}$ be a known $q$-dimensional linear subspace, $q\ge1$, with Frobenius-orthonormal basis $P_1,\ldots,P_q$. The unknown $A\in\R^{m\times n}$ is accessible through queries $x\mapsto Ax$ and $y\mapsto A^Ty$. One vector product counts as one query; a block product with $s$ columns counts as $s$ queries. All query directions below are fixed before observing any oracle response.

Define the basis-invariant positive-semidefinite matrix
\[
 R=\sum_{i=1}^q P_iP_i^T,\qquad \tr R=q.
\]
Fix a threshold $t>0$. Let $U$ have orthonormal columns spanning the eigenspaces of $R$ with eigenvalues greater than $t$. Write $\Pi=UU^T$ and $r=\rank\Pi$. Then $r\le q/t$. Query $A^TU$ and $Ag_1,\ldots,Ag_s$, where the $g_j$ are independent $N(0,I_n)$ vectors. These responses determine the objective
\begin{equation}\label{eq:objective}
 \ell(B)=\norm{U^T(A-B)}_F^2+
       \frac1s\sum_{j=1}^s\norm{(I-\Pi)(A-B)g_j}_2^2,
 \qquad B\in\mathcal L.
\end{equation}
Return a least-squares minimizer $\widetilde B$; a fixed minimum-coefficient-norm minimizer specifies the algorithm even when the sketched Gram matrix is singular.

\begin{theorem}\label{thm:main}
For any $\eps>0$, if
\begin{equation}\label{eq:samples}
 s\ge\max\left\{8(1+2t)\log(20q),\frac{160t}{\eps}\right\},
\end{equation}
then with probability at least $9/10$,
\begin{equation}\label{eq:relative}
 \norm{A-\widetilde B}_F^2
 \le(1+\eps)\min_{B\in\mathcal L}\norm{A-B}_F^2.
\end{equation}
In particular, writing $L_0=16\log(20q)+160/\eps$, the choice
\[
 t=\sqrt{q/L_0},\qquad
 s=\left\lceil8\log(20q)+tL_0\right\rceil
\]
uses at most
\begin{equation}\label{eq:universal}
 2\sqrt{qL_0}+8\log(20q)+1
\end{equation}
nonadaptive queries. The algorithm has polynomial-time dense preprocessing and least-squares postprocessing in an exact real-arithmetic model.
\end{theorem}

The parameter $\eps$ in this theorem is a squared-relative excess. A norm-relative guarantee follows with factor $\sqrt{1+\eps}\le1+\eps$. Conversely, a desired norm-relative factor $1+\varepsilon$ can be obtained by using squared parameter $2\varepsilon+\varepsilon^2$.

If the range of $R$ is measured exactly, no Gaussian queries are needed: all matrices in $\mathcal L$ vanish on its left orthogonal complement, whose contribution to the objective is independent of $B$. More generally one can optimize over every spectral split. If $\lambda_1\ge\cdots\ge\lambda_m\ge0$ are the eigenvalues of $R$, a split after $r$ eigenvectors has residual leverage $t=\lambda_{r+1}$ and query bound
\[
 r+\left\lceil\max\{8(1+2t)\log(20q),160t/\eps\}\right\rceil,
\]
with the separate zero-residual-leverage branch just described. The transposed construction based on $\sum_iP_i^TP_i$ is also available. The better of these plans can be selected using only the known family.

\section{The sketched Gram matrix}
Put
\[
 T_i=(I-\Pi)P_i,\qquad
 H_{i\ell}=\tr(T_i^TT_\ell),\qquad K=I_q-H.
\]
Both $H$ and $K$ are positive semidefinite, $H\preceq I_q$, and
\begin{equation}\label{eq:leverage}
 \sum_iT_iT_i^T=(I-\Pi)R(I-\Pi)\preceq tI_m.
\end{equation}
For a Gaussian vector $g$, let $M(g)$ be the $m\times q$ matrix with columns $T_i g$, and define
\[
 X(g)=K+M(g)^TM(g).
\]
Then $X(g)\succeq0$, $\E X(g)=I_q$, and the Gram matrix of \eqref{eq:objective} is
\[
 C=K+\frac1s\sum_{j=1}^sM(g_j)^TM(g_j)
   =\frac1s\sum_{j=1}^sX(g_j).
\]
The projected basis is \emph{not} reorthonormalized: keeping the exact identity $K+H=I_q$ is essential.

\begin{lemma}[Second moment]\label{lem:moment}
The random Gram contribution satisfies
\[
 \E(X(g)-I_q)^2\preceq2tI_q,
 \qquad \E X(g)^2\preceq(1+2t)I_q.
\]
\end{lemma}
\begin{proof}
For $a\in\R^q$, write $T=\sum_i a_iT_i$. Wick's formula for a standard real Gaussian vector gives
\begin{align*}
 a^T\E\bigl(M^TM-H\bigr)^2a
 =\sum_\ell\bigl\{
   \tr(T^TT_\ell T^TT_\ell)
   +\tr(T^TT_\ell T_\ell^TT)\bigr\}.
\end{align*}
For every square real matrix $Z$, $\tr Z^2\le\norm Z_F^2$. Hence each of the two terms inside braces is at most $\norm{T^TT_\ell}_F^2$. Using \eqref{eq:leverage},
\[
 a^T\E\bigl(M^TM-H\bigr)^2a
 \le2\tr\left(T^T\left(\sum_\ell T_\ell T_\ell^T\right)T\right)
 \le2t\norm T_F^2\le2t\norm a_2^2.
\]
Here the last inequality follows from $H\preceq I_q$. Since $X-I_q=M^TM-H$ and $\E(X-I_q)=0$, the second assertion follows from the first.
\end{proof}

\begin{lemma}[A matrix lower-tail estimate]\label{lem:tail}
Let $X_1,\ldots,X_s$ be independent positive-semidefinite $q\times q$ random matrices with $\E X_j=I_q$ and $\E X_j^2\preceq vI_q$. Then
\[
 \Prob\left\{\lambda_{\min}\left(\frac1s\sum_jX_j\right)\le\frac12\right\}
 \le q\exp\left(-\frac{s}{8v}\right).
\]
\end{lemma}
\begin{proof}
For $x\ge0$ and $\theta>0$, the scalar inequality
$e^{-\theta x}\le1-\theta x+\theta^2x^2/2$ gives by spectral calculus
\[
 \E e^{-\theta X_j}
 \preceq(1-\theta+\theta^2v/2)I_q
 \preceq e^{-\theta+\theta^2v/2}I_q.
\]
The Golden--Thompson trace inequality \cite{Golden1965} and conditioning yield the iteration
\begin{align*}
 \E\tr e^{-\theta\sum_{j=1}^sX_j}
 &\le e^{-\theta+\theta^2v/2}
       \E\tr e^{-\theta\sum_{j=1}^{s-1}X_j}\
 &\le q\exp\bigl(s(-\theta+\theta^2v/2)\bigr).
\end{align*}
To justify the first step, apply Golden--Thompson to the last summand and condition on all preceding summands. The conditional expectation of its exponential is bounded by a scalar multiple of the identity; multiplication inside a trace by a positive-semidefinite matrix preserves that bound. No commutativity of the $X_j$ is needed.

If $\lambda_{\min}(\sum_jX_j)\le s/2$, then the trace exponential is at least $e^{-\theta s/2}$. Markov's inequality therefore bounds the probability by
\[
 q\exp\bigl(s(-\theta/2+\theta^2v/2)\bigr).
\]
Choose $\theta=1/(2v)$. This is a standard matrix Laplace-transform argument in a particularly simple scalar-moment setting; see also \cite{Tropp2012}.
\end{proof}

Lemmas~\ref{lem:moment} and \ref{lem:tail} show
\begin{equation}\label{eq:gramgood}
 \Prob\{C\nsucceq I_q/2\}
 \le q\exp\left(-\frac{s}{8(1+2t)}\right).
\end{equation}

\section{Centered regression error}
Let $B_\star$ be the Frobenius projection of $A$ onto $\mathcal L$, and put $E=A-B_\star$. In particular $\ip{P_i}{E}_F=0$ for every $i$. In coefficient coordinates, the normal-equation offset at $B_\star$ is
\[
 z_i=\ip{U^TP_i}{U^TE}_F
       +\frac1s\sum_{j=1}^s\ip{T_i g_j}{(I-\Pi)E g_j}_2.
\]
Its expectation is zero because the deterministic and expected random parts sum to $\ip{P_i}{E}_F$.

\begin{lemma}\label{lem:noise}
The offset satisfies
\[
 \E\norm z_2^2\le\frac{2t}{s}\norm E_F^2.
\]
\end{lemma}
\begin{proof}
For a real square matrix $Z$ and $g\sim N(0,I)$,
\[
 \operatorname{Var}(g^TZg)=\tr Z^2+\tr(ZZ^T)\le2\norm Z_F^2.
\]
Independence of the $g_j$, followed by \eqref{eq:leverage}, gives
\begin{align*}
 \E\norm z_2^2
 &\le\frac2s\sum_i\norm{T_i^T(I-\Pi)E}_F^2\\
 &=\frac2s\tr\left(E^T(I-\Pi)
                  \left(\sum_iT_iT_i^T\right)(I-\Pi)E\right)\\
 &\le\frac{2t}{s}\norm E_F^2.
\end{align*}
\end{proof}

\begin{proof}[Proof of Theorem~\ref{thm:main}]
On the event $C\succeq I_q/2$, the coefficient error $d$ of $\widetilde B-B_\star$ is $C^{-1}z$, so $\norm d_2\le2\norm z_2$. If $\norm E_F>0$, Lemma~\ref{lem:noise} and Markov's inequality imply
\[
 \Prob\{\norm d_2^2>\eps\norm E_F^2,\ C\succeq I_q/2\}
 \le\frac{8t}{s\eps}.
\]
Together with \eqref{eq:gramgood}, the failure probability is at most
\[
 q\exp\left(-\frac{s}{8(1+2t)}\right)+\frac{8t}{s\eps}\le\frac1{10}
\]
under \eqref{eq:samples}. Pythagoras gives
$\norm{A-\widetilde B}_F^2=\norm E_F^2+\norm d_2^2$.
When $E=0$, the offset is identically zero, and recovery is exact on the same Gram event; no division by the optimal residual is needed.

For the stated universal choice, $s\ge8(1+2t)\log(20q)$ and $s\ge160t/\eps$. Also $r\le q/t$, whence
\[
 r+s\le q/t+tL_0+8\log(20q)+1
       =2\sqrt{qL_0}+8\log(20q)+1.
\]
The family orthonormalization, spectral decomposition, design formation, and least-squares solve use polynomially many real-arithmetic operations. This is not a bit-complexity or numerical backward-stability assertion.
\end{proof}

\section{Confidence amplification for affine families}
The same theorem applies to a known affine family after subtracting its known offset from each oracle response. For linear or affine families, run an odd number $L$ of independent copies using parameter $\eps/9$. From their coefficient vectors, choose one whose median Euclidean distance to all $L$ vectors is smallest, including its distance zero to itself.

\begin{proposition}\label{prop:amplification}
If $L$ is odd and $L\ge\log(1/\delta)/0.32$, the selected estimate satisfies \eqref{eq:relative} with probability at least $1-\delta$. All query plans may be selected in advance.
\end{proposition}
\begin{proof}
Each successful coefficient vector lies within
$\rho=\sqrt{\eps}\norm E_F/3$ of the true coefficient vector. If a strict majority are successful, every successful vector has median distance at most $2\rho$. The selected vector therefore has a ball of radius at most $2\rho$ containing a strict majority of all vectors. That ball intersects the successful majority, so the selected vector lies within $3\rho$ of the truth. Pythagoras gives \eqref{eq:relative}. Independent successes have probability at least $0.9$ each; Hoeffding's inequality bounds the probability of at most half successes by $\exp(-2(0.9-0.5)^2L)=\exp(-0.32L)$.
\end{proof}

The selector requires only finitely many pairwise distances. It does not estimate $\norm A_F$, the optimal residual, or unknown population losses.

\section{A closed-convex extension at constant success probability}
Let $\mathcal K$ be a nonempty closed convex subset of a known $q$-dimensional affine matrix space. Translate its known affine offset. Use the same queries. If the computed Gram matrix does not satisfy $C\succeq I_q/2$, report failure; otherwise minimize the sketched objective over $\mathcal K$. On this event the quadratic is coercive on the containing affine space, so it attains a unique minimum on the nonempty closed convex set. The reported failures are included in the probability bound below. This is a query-complexity statement conditional on carrying out that optimization; arbitrary descriptions of convex sets need not permit polynomial-time optimization.

\begin{proposition}\label{prop:convex}
The sample bound \eqref{eq:samples} gives, with probability at least $9/10$,
\[
 \norm{A-\widetilde B}_F^2\le(1+\eps)
                  \min_{B\in\mathcal K}\norm{A-B}_F^2.
\]
\end{proposition}
\begin{proof}
Let $B_\star$ now be the Frobenius projection onto $\mathcal K$, put $E=A-B_\star$, and write $d$ for the coefficient vector of $\widetilde B-B_\star$. Define $e_i=\ip{P_i}{E}_F$ and let $z$ be the sketched normal offset at $B_\star$ minus $e$. The centered calculation in Lemma~\ref{lem:noise} still gives $\E\norm z_2^2\le2t\norm E_F^2/s$; it did not require $E$ to be orthogonal to the containing affine space once the mean is subtracted.

Population projection optimality gives $e^Td\le0$. Sketched first-order optimality, tested against $B_\star$, gives
\[
 d^TCd\le(e+z)^Td.
\]
On $C\succeq I_q/2$ these imply $\norm d_2\le2\norm z_2$. The population excess obeys the stronger estimate
\begin{align*}
 \norm{A-\widetilde B}_F^2-\norm E_F^2
 &=\norm d_2^2-2e^Td\\
 &\le\norm d_2^2-2d^TCd+2z^Td\\
 &\le2z^Td\le4\norm z_2^2.
\end{align*}
The same Markov and Gram estimates prove the assertion. If $E=0$, the offset vanishes identically and the conclusion is exact on the Gram event.
\end{proof}

The median-radius amplification proposition is \emph{not} asserted for this constrained extension. For a general convex constraint, proximity to $B_\star$ alone does not upper-bound the population excess, because the normal term $-2e^Td$ need not vanish. A separate validation or confidence argument would be needed.

\section{Reproducibility and limitations}
The reference implementation uses only the two supplied oracle functions to access the target matrix. It exposes conservative proved plans separately from practical sample overrides, which carry no theorem-level guarantee. The verification programs check Gaussian fourth-moment identities, regression equations, exact-support recovery, zero-residual cases, the median-majority lemma, and coefficient-box constrained examples.

The mathematical analysis is in exact real arithmetic. The numerical implementation orthonormalizes a supplied basis using a floating-point rank tolerance and solves normal equations; it is not certified backward stable. The present note supplies an independently checkable argument, not an independently reviewed resolution or a certified priority claim.

\begin{thebibliography}{9}
\bibitem{Amsel2026}
N.~Amsel, P.~Avi, T.~Chen, F.~Duman Keles, C.~Hegde, C.~Musco, C.~Musco, and D.~Persson,
\emph{Query Efficient Structured Matrix Learning},
Proceedings of the Thirty Ninth Conference on Learning Theory,
PMLR \textbf{336}, 158--194 (2026).
\url{https://proceedings.mlr.press/v336/amsel26a.html}.
Preprint: \href{https://arxiv.org/abs/2507.19290}{arXiv:2507.19290}.

\bibitem{Prior2026}
Public repository account \texttt{andotheror},
\emph{Pure Relative Structured Matrix Learning with Square-Root Matvecs},
public manuscript announcement and repository README (2026).
\url{https://github.com/andotheror/matvec-structured-matrix-approx}.
The displayed deposit identifier is
\href{https://doi.org/10.5281/zenodo.21698599}{doi:10.5281/zenodo.21698599}.
The account name is not treated as verified personal authorship. Full-manuscript and deposit-metadata comparison remains outstanding.

\bibitem{Golden1965}
S.~Golden, \emph{Lower Bounds for the Helmholtz Function},
Physical Review \textbf{137}, B1127--B1128 (1965).
\href{https://doi.org/10.1103/PhysRev.137.B1127}{doi:10.1103/PhysRev.137.B1127}.

\bibitem{Tropp2012}
J.~A.~Tropp, \emph{User-Friendly Tail Bounds for Sums of Random Matrices},
Foundations of Computational Mathematics \textbf{12}, 389--434 (2012).
\href{https://doi.org/10.1007/s10208-011-9099-z}{doi:10.1007/s10208-011-9099-z}.
\end{thebibliography}
\end{document}

% END REVIEWED BODY
